\documentclass{homework}
% \usepackage{graphicx} % Required for inserting images
\usepackage{amsmath}
\usepackage{gensymb}
\usepackage{subfig}
\usepackage{float}
\usepackage{amsthm}

\def\one{\mbox{1\hspace{-4.25pt}\fontsize{12}{14.4}\selectfont\textrm{1}}} % 11pt    

\class{ECEN 5478}
\title{Project Check-in}
\author{Sean Jennings}
\date{December 7, 2023}

\begin{document}
\maketitle

\section{Context}

% Grid-level battery energy storage systems (BESSs) are projected to expand
% substantially by 2025. 



% With increased penetration of non-dispatchable renewable generation resources,
% energy storage solution will play a key role in the future of the power grid.
% Battery energy storage systems BESS are playing an increasingly large
% as the price of lithium ion batteries drops. In particular, a number
% of utility-scale storage projects have recently been installed or are
% in the project pipeline within the US, especially in states such as
% California with high renewable penetration.

% Fires at BESS facilities, such at the Elkhorn BESS (September 2022) in Moss Landing, CA,
% and the McMicken BESS (April 2019) in Surprise, AZ, suggest that the
% commercialization of grid-level BESS technology remains in a nascent state. 

% Within a deregulated energy market, 
% We will look at energy arbitrage oppurtunites for grid-level BESS plants within
% the deregulated wholesale energy market.

% Maximizing profitability of individual BESS plants is key to incentivizing
% the construction of additional energy storage.

% TODO:
% need to check how bidding works: model this??
% especially do we denote the suggested power capacity of the bid?

Due to the inherent variability of renewable generation resources, 
energy storage systems (ESSs) are expected to play an increasingly important role 
in the future power grid. Within a deregulated energy market, prices are lower
during periods of high renewable generation since the marginal cost of
renewable generation is relatively small compared to traditional fuel-based sources.
Conversely, during periods of low renewable production and high load, prices
are relatively elevated. This price differential represents an economic
oppurtunity for ESSs.

Market projections show that a large portion of the overall battery capacity in the
future will be stored in EVs. However, a variety of factors limit the dispatchability
of EVs as an ESS (e.g. EV owner driving patterns, the physical location of the EV
relative to EV charging infrastructure). Therefore, it is likely that utility-scale
stationary ESSs will represent an important asset in ensuring technical and
economic feasibility in a transition to power system with high penetrations of
renewable resources.

In particular, a number of grid-level battery energy storage (BESS) plants are
currently being (and planned to be) installed in the US, particularly in California and Texas.
The lifetime of batteries are limited by electrochemical degradation which limits
the economic utility which they provide.
For this project, we propose the consideration of the operational optimization of a BESS
in a energy time-shifting (arbitrage) application.

\section{Problem Formulation}

Consider an ESS unit which participates in a real-time energy market which is sufficiently
competitive that the bidding behavior of the unit has no effect upon price of energy
within the market (i.e. the unit can choose to purchase or sell energy at market price at will).
The market dispatches load/generation requests every $\Delta t$ minutes. We will
refer to the market dispatch action as a "market step." Suppose that at each market step $t$,
the unit has a perfect forcast of market prices per MWh, $p_t$, over the lifetime of the
$ESS$ unit, assumed to be $\bar{T}$ market steps, where $t,\bar{T}\in\N$, $p_t\in\R_{\geq0}^T$.

Due to computational constraints, it is clearly not possible to optimize over the
entire lifetime $T_L$, which in time is likely many years. Therefore, we select moving
time horizon $\mathcal{T}_t := [t, t + T]$, where $T\in\N$ is the "lookahead time", with $T<<T_L$,
and $t\in[0, T_L - T]$, is the time at which the prediction is made.

Let $\hat{u}_\tau\in\hat{\U}$ be a prediction of the power (in MW) discharged/charged by the unit
at times $\tau\in\mathcal{T}_t$. Suppose the power of the battery at each $\tau$
limited by $\U:=[\underbar{P}, \bar{P}]$,
so that $\hat{\U}$ is given by $\hat{\U} = \U \times (T \text{\it times}) \dots \times \U$, with
$\underbar{P}, \bar{P} \in \R$, $\underbar{P} < 0 < \bar{P}$.

Assume the battery degrades at a rate porportional to $\norm{u_\tau}^2$, so that
the monetary cost of a charge/discharge action (capital costs of the battery amortized over its
lifetime) is also porportional to $\norm{\hat{u}_t}^2$, with proportionality\footnote{
    A more realistic would be that the degradation rate is proportional to
    $\lVert u \rVert_1$; however, this results in a cost function that is neither smooth
    nor strongly convex.
} constant $\gamma$.
Let $\hat{\Phi}_t: \hat{\U} \to \R$ be the predicted cost of operation over the forcast horizon $T$
(negative values of $\hat{\Phi}_t$ indicate a net profit). Then, $\hat{\Phi}_t$ is given by
\begin{equation*}
    \hat{\Phi}_t(u) = \gamma \norm{u}^2 - p^T u / \Delta t.
\end{equation*}

Let $\X:=[\underbar{E}, \bar{E}]$ be the battery capacity limits with $\underbar{E}\in\R_{\geq0}$
indicating the minimum allowed capacity, and $\bar{E}\in\R_{\geq0}$ the maximum capacity,
$\underbar{E} < \bar{E}$.
Let $x_t\in\X$ be the state of charge (SoC) of the battery at time $t$. Assume the battery has
unit efficiency so that the discretized system dynamics are given by
\begin{equation*}
    x_{t+1} = x_t - \Delta t u_t
\end{equation*}
where $u_t\in\U$ is the battery power at time $t$, and is taken from the first element
of the control prediction $\hat{u}_t$ which is recomputed at each $t$.

We will solve the following time-varying optimization problem. Given SoC $x_t \in \X$,
\begin{align*}
    (P) \qquad &\min_{u \in \hat{U}} \hat{\Phi}_t(u) \\
        & \text{subj. to} \\
        & \qquad x_{\tau + 1} = x_\tau - \Delta t u_\tau \\
        & \qquad x_\tau \in \X \text{ for each } \tau \in \mathcal{T}_t
\end{align*}
where $u_\tau$ is the $(\tau - t)$th entry of $u$. We note that this problem represents
merely an approximation of the global problem $P_g$. Given $x_0 \in \X$ minimize:
\begin{align*}
    (P_g) \qquad &\min_{u \in \hat{U}} \hat{\Phi}_t(u) \\
        & \text{subj. to} \\
        & \qquad x_{t + 1} = x_t - \Delta t u_t \\
        & \qquad x_t \in \X \text{ for each } t \in [0, T_L]
\end{align*}

Let $u^* \in \R^{T_L}$ be the solution to $P_g$ and $\hat{u}_t \in \R^T$ be the solution
to $P$. Then, since we select just the first element, $\hat{u}_{t,0}$, before making a new
prediction at $t+1$, the distance from the optimal solution, $e_t$, is given by:
\begin{equation*}
    e_t := \norm{\hat{u}_{t,0} - u_t^*}
\end{equation*}
% We propose the following as potential extensions to increase the challenge of the problem:
% \begin{enumerate}
%     \item Additionally providing frequency control services
%         (need to decide what timescale to model freq control on)
%     \item precisely model degradation costs (violates strong convexity/smoothness conditions)
%     \item Representation of distributed control aspects (plant-level vs module-level BMS)
%     \item Modeling of low-level LiPo mechanics
% \end{enumerate}

\subsection{Constraint Reformulation}

It is not immediately clear whether the constraints $x_\tau\in\X$, within $P$ and $P_g$ represent
a convex set. We reformulate constraints into a new set, $\hat{U}_{t,\X}\subseteq \hat{U} \subset \R^T$,
with dependence on $\X$, which accounts for SoC-imposed constraints and
will demonstrate the convexity of problems $P$ and $P_g$.

Consider the more general case of an discrete LTI system given by:
\begin{equation}
    x_{t + 1} = A x_t + B u_t
    \label{eq:lti-system}
\end{equation}
where $x_t\in \R^{N_x}$, $u_t\in\R^{N_u}$ for each $t$, $A\in\R^{N_x \times N_x}$
$B \in \R^{N_x \times N_u}$, and $N_x,N_u\in\N$ are the state and control dimensions
repectively.

Then, given a control $u_\tau$, for each $\tau \in [t, \dots, t + T - 1]$, we have
\begin{equation}
    x_{\tau + 1} = A^{\tau + 1 - t} x_t + \sum_{i=t}^{\tau} A^{\tau - i} B u_i
    \label{eq:integrator}
\end{equation}

\begin{proof}
    We prove by induction. The base case of $\tau = t$ holds since
    (\ref{eq:integrator}) reduces to (\ref{eq:lti-system}).

    Suppose that (\ref{eq:integrator}) holds for $\tau$, then
    \begin{align*}
        x_{\tau + 2} &= A x_{\tau + 1} + Bu_{\tau + 1} \\
            &= A \biggl(
            A^{\tau + 1 - t} x_t + \sum_{i=t}^\tau A^{\tau - i} B u_i
        \biggr) + B u_{\tau + 1} \\
        &= A^{\tau + 2 - t} x_t + \biggl(\sum_{i=t}^\tau A^{(\tau + 1) - i} B u_i\biggr) + B u_{\tau + 1}\\
        &= A^{\tau + 2 - t} x_t + \sum_{i=t}^{\tau + 1} A^{(\tau + 1) - i} B u_i.
    \end{align*}
    So (\ref{eq:integrator}) also holds for $\tau + 1$.
\end{proof}

We rewrite (\ref{eq:integrator}) in matrix form:
\begin{equation*}
    \hat{x} := \mat{
        x_{t + 1} \\
        x_{t + 2} \\
        \vdots \\
        x_{t + T}
    } = \mat{
        B & 0 & \cdots & 0 \\
        AB & B & \cdots & 0 \\
        \vdots & \vdots & \ddots & \vdots \\
        A^{T - 1} & A^{T - 2} & \cdots & B
    } \mat{
        u_t \\ u_{t+1} \\ \vdots \\ u_{T - 1}
    } + \mat{
        A \\ A^2 \\ \vdots \\ A^T
    } x_t
\end{equation*}

Putting $A = 1$ and $B = - \Delta t$, we have
\begin{equation*}
    \hat{x} = \Gamma \hat{u} + \one x_t
\end{equation*}


where $\hat{x}$ are the predicted states from $t+1$ to $T+1$,
$\Gamma \in \R^{T \times T}$ is lower triangular matrix with all non-zero entries
equal to $-\Delta t$, and $\one$ is the vector of ones. 

Since we have the following constraints on $\hat{x}$,
\begin{equation*}
    \one \underbar{E} \leq \hat{x} \leq \one \bar{E},
\end{equation*}

we can define $\hat{U}_{t, \X}$, the SoC-imposed constraints at time $t$, as
\begin{equation*}
    \hat{U}_{t,\X} = \biggl\{
        \hat{u} \in \R^T: \mat{\Gamma \\ -\Gamma} \hat{u} \leq
            \mat{\one (\bar{E} - x_t) \\ \one (x_t - \underbar{E})}
    \biggr\}.
\end{equation*}
This set is a convex polytope, so
the optimization problem is convex ($\Phi_t$ is $\gamma$-strongly convex and $\gamma$-smooth).
We can then solve the alternative formulation $P'$:

\begin{align*}
    (P') \qquad &\min_{u \in \hat{U}_{t,\X}} \hat{\Phi}_t(u) \\
\end{align*}

\section{Sketch of Algorithm}
We will use the projected gradient descent method to track the time-varying solution of $P$.

\begin{equation}
    \hat{u}_{t+1} = \text{proj}_{\hat{U}_{t,\X}} \biggl\{
        \hat{u}_t - \alpha \Phi_t(\hat{u}_t)
    \biggr\}
\end{equation}

with $\alpha\in\R$ as the step-size of the gradient method.

However, since the constraints on each entry $u_\tau$ of $\hat{u}$ are not independent
of other entries, the author is not aware of a closed form of the projection
function $\text{proj}_{\hat{\U}'}$ for a general convex polytope. 
We will need to solve the optimization subproblem
\begin{equation*}
    (P_{sub}) \qquad \text{proj}_{\hat{\U}'}(u) = \text{arg}\min_{v\in\hat{\U}'}
        \norm{u - v}
\end{equation*}
at each timestep.

\section{Next Steps}

For the next steps on this project, we will:
\begin{enumerate}
    \item Derive a bound on $e_t$ as a function of the time horizon $T$.
    \item Implement the above algorithm in Python.
    \item Pull LMP data from CAISO and interpolate between data points to simulate a faster time-varying
        problem than the 5-min steps sizes.
\end{enumerate}

\end{document}
